% ===================================================================
%  اللوحة رقم ١٣ — الفندق. — المطعم. — المقهى.
%  Traduction 2026 d'apres texto/io, controlee sur texto/fr, texto/en
%  et texto/es. Consigne : voir texto/ar/10-lawha-01.tex.
%
%  Deux notes, marquees toutes deux « (*) », et deux appels : celui du
%  bloc 01-2 (la chambre) et celui du bloc 09-1 (le menu). L'ordre les
%  departage.
%
%  CAEN S'ECRIT كاين ET NON كان. Le nom de la ville se prononce en un
%  seul souffle -- « kan » --, et كان seul est le verbe « etre » au
%  passe, le mot le plus ordinaire de la langue : « كرشًا على طريقة
%  كان » se lisait « des tripes a la maniere dont c'etait ». Le ya
%  leve l'homographie sans changer le son, et c'est la graphie recue
%  pour cette ville.
% ===================================================================

%%K t13-noto-1 noto
\VUnotes{143.2mm}{%
\textit{\VUgras{(*) لتفاصيل غرفة النوم انظر اللوحة رقم ٦.}}%
}

%%K t13-apar-1 apar
\VUsaut{3.0mm}
\VUtitre{15.0pt}{60}{اللوحة رقم ١٣}
\VUsaut{1.6mm}
\VUpk{10.2pt}{40}{الفندق. --- المطعم.}
\VUsaut{2.0mm}
\VUpk{10.2pt}{40}{المقهى.}
\VUsaut{3.4mm}

%%K t13-01-1 p
1. --- لمّا وصلت إلى رويان مساء أمس، نزلت في \textit{فندق المحيط} الذي
أوصاني به صديق.

%%K t13-01-2 p
وأعطوني \VUgras{غرفة} \textsuperscript{(2)} جميلة في \VUgras{الطابق}
\textsuperscript{(1)} الثاني، فنمت فيها نومًا طيبًا جدًا (*).

%%K t13-02-1 p
2. --- وهذا الصباح، وأنا خارج من غرفتي حيث جاءتني \VUgras{الخادمة}
\textsuperscript{(3)} بالفطور، دخلت \VUgras{غرفة التدخين} \textsuperscript{(5)}، وهي
أيضًا \VUgras{غرفة المطالعة} \textsuperscript{(4)}، لأتصفح \VUgras{الجرائد}
\textsuperscript{(6)} وأنا أدخّن \VUgras{سيجارة} \textsuperscript{(8)}. ولما فرغت من
القراءة نزلت بـ\VUgras{المصعد} \textsuperscript{(9)} وخرجت أتنزه في المدينة.

%%K t13-03-1 p
3. --- ولأني نسيت أن أسأل \VUgras{موظف} \textsuperscript{(12)} الفندق عن ساعة
تقديم الطعام على \VUgras{المائدة المشتركة} \textsuperscript{(11)}، عدت مبكرًا
واغتنمت الفرصة لأستحم في \VUgras{حمّام} \textsuperscript{(10)} الفندق. ثم ذهبت
إلى \VUgras{الصندوق} \textsuperscript{(15)} لأتصل بأحد أصدقائي هاتفيًا. لكن
\VUgras{الهاتف} \textsuperscript{(13)} كان مشغولًا؛ ولما رأت \VUgras{أمينة
الصندوق} \textsuperscript{(14)} حيرتي --- وكانت تقيّد \VUgras{فاتورة}
\textsuperscript{(17)} في \VUgras{سجل النزلاء} \textsuperscript{(16)} --- طلبت من
\VUgras{البوّاب} \textsuperscript{(18)} أن يرسل \VUgras{الصبي} \textsuperscript{(19)}
ليقضي حاجتي على \VUgras{دراجته} \textsuperscript{(20)}.

%%K t13-04-1 p
4. --- فشكرتها وخرجت لأعطي إكرامية لـ\VUgras{الحمّال} \textsuperscript{(24)}
(الرجل لكل شيء) الذي جلب لي من المحطة على \VUgras{عربته اليدوية}
\textsuperscript{(25)} \VUgras{علبة قبعتي} \textsuperscript{(26)}
و\VUgras{حقيبتي} \textsuperscript{(27)} و\VUgras{كيس سفري} \textsuperscript{(28)}. وناولني
\VUgras{ساعي البريد} \textsuperscript{(21)}، وقد وصل في تلك اللحظة،
\VUgras{رسالة} \textsuperscript{(22)}.

%%K t13-05-1 p
5. --- وبقيت قليلًا بعدُ على العتبة، قرب \VUgras{صندوق البريد}
\textsuperscript{(23)}، أرقب حركة الشارع. وكان \VUgras{سائق} \textsuperscript{(30)}
\VUgras{العربة العامة} \textsuperscript{(29)} يحمّل على مركبته \VUgras{صندوق}
\textsuperscript{(31)} أحد \VUgras{النزلاء} \textsuperscript{(32)} الذي كان \VUgras{صاحب
الفندق} \textsuperscript{(33)} يودّعه. وكان \VUgras{عامل برق} \textsuperscript{(35)} شاب
يحمل، غير مستعجل البتة، \VUgras{برقية} \textsuperscript{(34)} لأحد نزلاء الفندق؛
و\VUgras{سائح} \textsuperscript{(36)} و\VUgras{آلة تصويره} \textsuperscript{(38)} على
كتفه يراجع \VUgras{دليله} \textsuperscript{(37)}.

%%K t13-tit-1 sub
\VUsaut{4.0mm}
\VUpk{10.2pt}{40}{ساعة الغداء.}
\VUsaut{3.0mm}

%%K t13-06-1 p
6. --- أمام السياج كان \VUgras{رجل مشاوير} \textsuperscript{(39)} وادع يغفو بين
\VUgras{كلّابيه} \textsuperscript{(40)} و\VUgras{علبة تلميع الأحذية}
\textsuperscript{(41)}، بينما كانت \VUgras{غسّالة} \textsuperscript{(42)} شابة تحمل
\VUgras{سلة} \textsuperscript{(43)} من الثياب وتضحك من وقار \VUgras{رئيس الخدم}
\textsuperscript{(44)} الواقف عند الباب.

%%K t13-07-1 p
7. --- ومنظر \VUgras{الطاهي} \textsuperscript{(45)}، وقد خرج من \VUgras{مطبخه}
\textsuperscript{(47)} في \VUgras{القبو} \textsuperscript{(48)} ليرسل \VUgras{صبي
مطبخ} \textsuperscript{(46)} في حاجة، ذكّرني بأن ساعة الغداء قد قربت، فذهبت إلى
المطعم في الجهة الأخرى من الفناء، حيث كان \VUgras{سائق سيارة}
\textsuperscript{(50)} يركن \VUgras{سيارته} \textsuperscript{(49)}، بينما كان
\VUgras{ميكانيكي} \textsuperscript{(52)} يصلح، على إرشاد \VUgras{راكب الدراجة}
\textsuperscript{(53)}، \VUgras{دراجة نارية ثلاثية} \textsuperscript{(51)} أصابها حادث
لتوّها.

%%K t13-08-1 p
8. --- وسألت \VUgras{خادمًا} \textsuperscript{(54)} شابًا يحمل \VUgras{أقداح جعة}
\textsuperscript{(55)} أكانت المائدة المشتركة تحت الشرفة، فلما قال نعم، جلست
إلى إحدى الموائد وقُدّم لي طعام ممتاز.

%%K t13-noto-2 noto
\VUnotes{143.2mm}{%
\textit{\VUgras{(*) بمعونة قائمة الأطباق التي يعطيها المعجم، يستطيع
المعلم أن ينوّع بسهولة قائمة الطعام الآتية.}}%
}

%%K t13-tit-2 sub
\VUsaut{4.0mm}
\VUpk{10.2pt}{40}{غداء ممتاز.}
\VUsaut{3.0mm}

%%K t13-09-1 p
9. --- جاءني النادل بقائمة الغداء (*) وقائمة الخمور. فاخترت وطلبت
\VUgras{قريدسًا} بـ\VUgras{الزبدة} \VUgras{مقبّلات}، وفي الطبق الأول
\VUgras{كرشًا على طريقة كاين}. ولم آخذ \VUgras{سمكًا}، لكني طلبت حصة من
\VUgras{فخذ} الخروف، ومن \VUgras{الخضر} \VUgras{سبانخ} بالقشدة. ثم، إذ لم
يرقني شيء من \VUgras{الأطباق الحلوة}، أمرت بإحضار أصناف من الحلوى:
\VUgras{جبن} و\VUgras{فاكهة} و\VUgras{بسكويت}. وأضفت \VUgras{خمرًا} ممتازة
من سان-إميليون، وقمت عن المائدة راضيًا جدًا عن طعامي بعد أن أعطيت
\VUgras{النادل} \textsuperscript{(57)} \VUgras{إكرامية} \textsuperscript{(59)} سخية.

%%K t13-10-1 p
10. --- ولما رآني \VUgras{المدير} \textsuperscript{(60)} مستعدًا للانصراف، اقترح
عليّ أن أخرج إلى الشرفة لأتأمل منظر \VUgras{المحيط} \textsuperscript{(62)} البديع.
وفي البعد كانت تُتبيَّن \VUgras{قوارب} \textsuperscript{(63)} صيد صغيرة
و\VUgras{سفن شراعية} \textsuperscript{(64)} و\VUgras{باخرة} \textsuperscript{(65)} هائلة
يرتسم شبحها الأسود على \VUgras{الغيوم} \textsuperscript{(67)} البيضاء عند
\VUgras{الأفق} \textsuperscript{(66)}. ونسيم خفيف يحرّك \VUgras{العَلَم}
\textsuperscript{(70)} المغروز فوق \VUgras{لافتة} \textsuperscript{(69)}
\VUgras{البهو} \textsuperscript{(68)}.

%%K t13-tit-3 sub
\VUsaut{3.0mm}
\VUpk{10.2pt}{40}{التسلية.}
\VUsaut{3.0mm}

%%K t13-11-1 p
11. --- وكان المنظر من الطرافة بحيث عزمت أن أصعد في الغد إلى شرفة المقهى
ومعي \VUgras{منظار مسرح} \textsuperscript{(71)} أو \VUgras{منظار مقرّب}
\textsuperscript{(72)}.

%%K t13-12-1 p
12. --- ولما تركت الشرفة ذهبت إلى مقهى الفندق لأتناول \VUgras{مشروبًا}
\textsuperscript{(73)} وألعب دورًا من \VUgras{البلياردو} \textsuperscript{(75)}. وعبرت
قاعة المقهى حيث كان كثير من الزبائن يلعبون \VUgras{الشطرنج} \textsuperscript{(78)}
و\VUgras{الداما} \textsuperscript{(79)} و\VUgras{الدومينو} \textsuperscript{(80)}
و\VUgras{الطاولة} \textsuperscript{(81)} أو \VUgras{الورق} \textsuperscript{(82)}،
وصعدت رأسًا إلى \VUgras{قاعة البلياردو} \textsuperscript{(74)}.

%%K t13-13-1 p
13. --- ولما انتهى الدور نزلت إلى الطابق الأرضي وطلبت من النادل
\VUgras{ظرفًا} \textsuperscript{(84)} و\VUgras{ورقًا} \textsuperscript{(83)}
و\VUgras{طابع بريد} \textsuperscript{(85)} و\VUgras{قلمًا} \textsuperscript{(87)}
و\VUgras{حبرًا} \textsuperscript{(86)} لأكتب رسالة.

%%K t13-13-2 p
ثم ناديت \VUgras{حوذيًا} \textsuperscript{(92)} كانت \VUgras{عربته}
\textsuperscript{(88)} واقفة أمام المقهى. وسألته عن أجره بالساعة وبالمشوار، فلما
اتفقنا فتح لي \VUgras{باب العربة} \textsuperscript{(89)} وأجلسني في الداخل. ثم
عاد إلى \VUgras{مقعده} \textsuperscript{(91)}، وحثّ الخيل حثًّا نشيطًا
بـ\VUgras{سوطه} \textsuperscript{(93)}، وانطلقنا في نزهة بديعة.
\VUsaut{6.0mm}
\VUfilet{22mm}
